\documentclass[10pt,a4paper]{amsart}
\usepackage[margin=19mm]{geometry}
\usepackage{amsmath,amssymb,mathtools}
\usepackage[scr=rsfs,cal=euler]{mathalfa}
\usepackage{xfrac}
\usepackage[unicode,hidelinks,bookmarks=false]{hyperref}
\newcommand{\ie}{i.e.\ }
\newcommand{\cf}{cf.\ }
\newcommand{\resp}{respectively\ }
\newcommand{\Subsection}{\subsection}
\providecommand{\nobreakdash}{\nobreak\hbox{-}}
\setlength{\emergencystretch}{2em}
\multlinegap0pt
\usepackage{bm}
\usepackage{bbm}
\usepackage{dsfont}
\usepackage{xspace}
\usepackage{cancel}
\usepackage{mathtools}
\usepackage{amsthm}
\makeatletter
\@ifundefined{theorem}{\newtheorem{theorem}{Theorem}}{}
\@ifundefined{lemma}{\newtheorem{lemma}[theorem]{Lemma}}{}
\makeatother
\title[Topological Dynamics of Pullback Maps on Full Shifts]{Topological Dynamics of Pullback Maps on Full Shifts}
\author{Alonso Castillo-Ramirez, Luguis De Los Santos Ba\~nos}
\date{Source: arXiv:2607.06423v1 (CC BY 4.0)}
\begin{document}
\maketitle

\noindent\emph{Source and licence.} Topological Dynamics of Pullback Maps on Full Shifts. Version arXiv:2607.06423v1 (CC BY 4.0), distributed under the Creative Commons Attribution 4.0 International licence, as recorded in the arXiv licence metadata for this exact version and shown on the abstract page https://arxiv.org/abs/2607.06423v1. Full rights notice: licenses/arxiv-2607.06423-CC-BY-4.0.txt.\par\smallskip

\noindent\textit{Editorial scope.} Counted unit: the Theorem whose source label is \texttt{thm:general-mixing} in the article (source0.tex lines 435--437, source label \texttt{thm:general-mixing}), delivered together with the article's own complete original proof of it (source0.tex lines 439--485, one \texttt{proof} environment of 47 lines). No mathematics is written, repaired or re-derived: statement and proof are the article's own text line for line. Required context retained uncounted: lem:finite-escape (lines 417--427). Every definition, display and label of those lines is retained unchanged. Omitted from this package and not redistributed: 1--416, 428--434, 438--438, 486--644. Nothing inside a retained excerpt is omitted or altered: every retained body is delivered exactly as the article's lines are written, so no omission receipt of any kind accompanies this package. Citations: the retained excerpts carry no citation command at all, so no bibliography entry of the article is transcribed and no reference is redistributed. Reference adaptation: none was needed and none was made; every reference inside the delivered unit resolves inside it, so no unresolved reference or citation is produced. Printed slips kept literal: no printed word, symbol, hypothesis or number of the counted statement or of its proof was altered. Layout and class: the article's own class is \texttt{article} with packages including inputenc, fontenc, amsmath, amssymb, amsthm, geometry, xcolor, authblk, hyperref; the wrapper is the lane's pinned \texttt{amsart} editorial wrapper at 10pt on A4, which restates the theorem-like environments the retained text needs and adds the article's own macro definitions that the retained text needs (none), with extra wrapper packages (bm, bbm, dsfont, xspace, cancel, mathtools). The delivered numbering is the unit's own. Assets: no retained span contains a figure environment, a graphics-inclusion command, a PDF-page inclusion, a table or a TikZ picture; no image file is copied into this package, the package declares no asset dependency and no image file is redistributed with it. Licence: version arXiv:2607.06423v1 of this article is distributed under the Creative Commons Attribution 4.0 International licence, as recorded in the arXiv licence metadata for this exact version and shown on the abstract page https://arxiv.org/abs/2607.06423v1. A full rights notice is delivered at licenses/arxiv-2607.06423-CC-BY-4.0.txt. Characters: every retained excerpt is pure ASCII, so the delivered engine, XeLaTeX with its default Latin Modern text font, reports no missing character.
\begin{lemma}\label{lem:finite-escape}
Let $\phi:G\to G$ be an endomorphism. The following are equivalent:
\begin{enumerate}
\item No nontrivial element of $G$ is eventually periodic under $\phi$.
\item For every finite set $F\subseteq G\setminus\{e_G\}$ and every finite set $E\subseteq G$, there exists $N\geq 0$ such that
\[
\phi^n(F)\cap E=\emptyset
\qquad\text{for all }n\geq N.
\]
\end{enumerate}
\end{lemma}

\begin{theorem}\label{thm:general-mixing}
Let $\phi:G\to G$ be an endomorphism. For each $c\in A$, the system $(X_c,\phi^*)$ is topologically mixing if and only if no nontrivial element of $G$ is eventually periodic under $\phi$.
\end{theorem}

\begin{proof}
Suppose first that there exists a nontrivial eventually periodic element $g\in G$.

If the eventual cycle of $g$ contains a nontrivial element, then there exist $k\geq 0$ and $p\geq 1$ such that $h=\phi^k(g)\neq e_G$ and $\phi^p(h)=h$. Choose distinct symbols $a,b\in A$ and consider
\[
U=\{x\in X_c:x(h)=a\},
\qquad
V=\{x\in X_c:x(h)=b\}.
\]
For every $x\in X_c$ and every $m\geq 1$,
\[
(\phi^*)^{mp}(x)(h)=x(\phi^{mp}(h))=x(h).
\]
Therefore $(\phi^*)^{mp}(U)\cap V=\emptyset$ for all $m\geq 1$, and the system is not mixing.

If the orbit of $g$ eventually reaches $e_G$, then there exists $N\geq 0$ such that $\phi^n(g)=e_G$ for all $n\geq N$. Choose $d\in A\setminus\{c\}$ and set
\[
V=\{x\in X_c:x(g)=d\}.
\]
For every $x\in X_c$ and every $n\geq N$,
\[
(\phi^*)^n(x)(g)=x(e_G)=c.
\]
Hence $(\phi^*)^n(V)\cap V=\emptyset$ for all $n\geq N$, so the system is not mixing.

Conversely, assume that no nontrivial element is eventually periodic. Let $U,V\subseteq X_c$ be nonempty open sets. Choose cylinders
\[
C(F_1,q_1)\cap X_c\subseteq U,
\qquad
C(F_2,q_2)\cap X_c\subseteq V,
\]
where, after removing the identity coordinate if necessary, we may assume $F_1,F_2\subseteq G\setminus\{e_G\}$. By Lemma~\ref{lem:finite-escape}, there exists $N\geq 0$ such that
\[
F_1\cap \phi^n(F_2)=\emptyset
\qquad\text{for all }n\geq N.
\]
Moreover, $\phi^n(F_2)$ does not contain $e_G$, since otherwise an element of $F_2$ would eventually reach $e_G$ and hence would be eventually periodic. For each $n\geq N$, define $x_n\in A^G$ by imposing
\[
x_n(e_G)=c,
\qquad
x_n|_{F_1}=q_1,
\qquad
x_n(\phi^n(g))=q_2(g)\quad (g\in F_2),
\]
and extending arbitrarily elsewhere. These conditions are compatible because the hypothesis of the theorem also implies that $\phi$ is injective: otherwise a nontrivial element
of $\ker(\phi)$ would be eventually periodic. Therefore every iterate $\phi^n$ is also injective. Then $x_n\in U$ and $(\phi^*)^n(x_n)\in V$. Therefore $(\phi^*)^n(U)\cap V\neq\emptyset$ for all $n\geq N$, proving mixing.
\end{proof}

\end{document}
